\chapter{Het scalair product in het vlak}\label{ch:g11:scal}

Het \omterm{def:g11:scal:dot}{scalair product} hangt een \emph{getal} aan een paar \omterm{def:g11:vect:vector}{vectoren}, en dat getal
ziet meetkunde: het verdwijnt precies bij loodrechte richtingen, en het meet
hoeken en lengten. Het is de motor achter de cosinusregel, achter
\omterm{def:g10:algebra:equation}{vergelijkingen} van cirkels en achter berekeningen met loodrechte stand.
Hetzelfde gereedschap, in de ruimte, is het onderwerp van \cref{ch:g12:space}.

\section{Definitie en eerste berekeningen}

\begin{definition}[Scalair product]\label{def:g11:scal:dot}
Het \emph{scalair product}\index{scalair product} van twee \omterm{def:g11:vect:vector}{vectoren} $\vec u$
en $\vec v$ is het getal
\[
\vec u \cdot \vec v
= \tfrac12\left(\norm{\vec u + \vec v}^2 - \norm{\vec u}^2
- \norm{\vec v}^2\right),
\]
waarbij $\norm{\vec u}$ de lengte (de \emph{norm}\index{norm}) van $\vec u$
aanduidt.
\end{definition}

\begin{theorem}[De vier gezichten van het scalair product]\label{thm:g11:scal:faces}
Zij $\vec u\,(x, y)$ en $\vec v\,(x', y')$ \omterm{def:g11:vect:vector}{vectoren} verschillend van
$\vec 0$, en $\theta$ de hoek ertussen. Dan geldt
\begin{enumerate}
  \item \emph{(\omterm{def:g10:coordgeom:system}{coördinaten})} $\vec u \cdot \vec v = xx' + yy'$;
  \item \emph{(\omterm{def:g11:trigo:cossin}{cosinus})} $\vec u \cdot \vec v =
        \norm{\vec u}\,\norm{\vec v}\cos\theta$;
  \item \emph{(projectie)} $\vec u \cdot \vec v = \vec u \cdot \vec v'$,
        waarbij $\vec v'$ de \omterm{def:g11:scal:orthogonal}{orthogonale} projectie van $\vec v$ op de
        richting van $\vec u$ is; is $\vec u = \vect{OA}$ en
        $\vec v = \vect{OB}$, dan is
        $\vec u \cdot \vec v = \overline{OA} \times \overline{OH}$, een
        product van getekende lengten, waarbij $H$ de voet is van de
        loodlijn uit $B$ op de rechte $(OA)$;
  \item \emph{(\omterm{def:g11:scal:dot}{normen})} de definiërende formule van \cref{def:g11:scal:dot}.
\end{enumerate}
\end{theorem}

\begin{proof}[Bewijs van 1 en 2]
\emph{1.} In \omterm{def:g10:coordgeom:system}{coördinaten} is $\norm{\vec u}^2 = x^2 + y^2$ (Pythagoras) en
heeft $\vec u + \vec v$ \omterm{def:g10:coordgeom:system}{coördinaten} $(x + x', y + y')$, dus
\begin{align*}
2\,\vec u \cdot \vec v
&= (x + x')^2 + (y + y')^2 - (x^2 + y^2) - (x'^2 + y'^2)\\
&= 2xx' + 2yy',
\end{align*}
na de kwadraten uit te werken en weg te strepen.

\emph{2.} Kies het \omterm{def:g10:coordgeom:system}{assenstelsel} zo dat $\vec u$ langs de positieve $x$-as
wijst: $\vec u\,(\norm{\vec u}, 0)$ en
$\vec v\,(\norm{\vec v}\cos\theta,\ \norm{\vec v}\sin\theta)$
(\cref{ch:g11:trigo}). Formule 1 geeft dan
$\vec u \cdot \vec v = \norm{\vec u}\,\norm{\vec v}\cos\theta + 0$. Formule
3 leest dezelfde berekening af: de projectie van $\vec v$ op de $x$-as heeft
getekende lengte $\norm{\vec v}\cos\theta = \overline{OH}$.
\end{proof}

\begin{omfigure}
\begin{tikzpicture}[scale=1.25]
  \coordinate (O) at (0,0);
  \coordinate (A) at (3.2,0);
  \coordinate (B) at (55:2.4);
  \coordinate (H) at (1.3767,0);
  \draw[-Stealth, omDef, thick] (O) -- (A)
    node[below, font=\small, pos=0.85] {$\vec u$};
  \draw[-Stealth, omThm, thick] (O) -- (B)
    node[above left, font=\small, pos=0.8] {$\vec v$};
  \draw[black, dashed, thin] (B) -- (H);
  \draw[thin] (H) ++(-0.18,0) -- ++(0,0.18) -- ++(0.18,0);
  \draw[omProp, very thick] (O) -- (H);
  \fill (H) circle (1.2pt) node[below, font=\small] {$H$};
  \node[below left, font=\small] at (O) {$O$};
  \draw[->,omThm,thin] (0.55,0) arc (0:55:0.55);
  \node[omThm, font=\small] at (27:0.78) {$\theta$};
\end{tikzpicture}

{\small Het \omterm{def:g10:functions:function}{beeld} van de projectie:
$\vec u \cdot \vec v = \overline{OA} \times \overline{OH}$, waarbij $H$ de
voet is van de loodlijn uit de punt van $\vec v$. Het oranje lijnstuk $OH$
heeft getekende lengte $\norm{\vec v}\cos\theta$.}
\end{omfigure}

\begin{example}\label{ex:g11:scal:compute}
$\vec u\,(3, 1)$ en $\vec v\,(2, -4)$: $\vec u \cdot \vec v = 6 - 4 = 2$. De
\omterm{def:g11:scal:dot}{normen} zijn $\norm{\vec u} = \sqrt{10}$ en $\norm{\vec v} = \sqrt{20}$, dus
$\cos\theta = \frac{2}{\sqrt{10}\sqrt{20}} = \frac{2}{10\sqrt2}
= \frac{\sqrt2}{10}$: de hoek is iets minder dan $82^\circ$.
\end{example}

\begin{proposition}[Rekenregels]\label{prop:g11:scal:rules}
Voor alle \omterm{def:g11:vect:vector}{vectoren} en alle $\lambda \in \R$ geldt
\[
\vec u \cdot \vec v = \vec v \cdot \vec u, \qquad
\vec u \cdot (\vec v + \vec w) = \vec u \cdot \vec v + \vec u \cdot
\vec w, \qquad
(\lambda \vec u)\cdot \vec v = \lambda\,(\vec u \cdot \vec v),
\]
\[
\vec u \cdot \vec u = \norm{\vec u}^2,
\qquad
\norm{\vec u + \vec v}^2 = \norm{\vec u}^2 + 2\,\vec u\cdot\vec v +
\norm{\vec v}^2 .
\]
\end{proposition}

\begin{proof}
Alles volgt uit de formule in \omterm{def:g10:coordgeom:system}{coördinaten}: zo is
$\vec u \cdot (\vec v + \vec w) = x(x' + x'') + y(y' + y'') =
(xx' + yy') + (xx'' + yy'')$. De laatste identiteit werkt
$\norm{\vec u + \vec v}^2 = (\vec u + \vec v)\cdot(\vec u + \vec v)$ uit met
de eerste drie regels --- het \omterm{def:g11:scal:dot}{scalair product} gedraagt zich precies als een
gewoon product, zodat de gebruikelijke merkwaardige producten ook op
\omterm{def:g11:vect:vector}{vectoren} van toepassing zijn.
\end{proof}

\section{Loodrechte stand}

\begin{definition}[Orthogonale vectoren]\label{def:g11:scal:orthogonal}
Twee \omterm{def:g11:vect:vector}{vectoren} heten \emph{orthogonaal}\index{orthogonale vectoren},
geschreven $\vec u \perp \vec v$, wanneer $\vec u \cdot \vec v = 0$.
\end{definition}

\begin{remark}
Voor \omterm{def:g11:vect:vector}{vectoren} verschillend van $\vec 0$ toont formule 2 van
\cref{thm:g11:scal:faces} dat $\vec u \cdot \vec v = 0 \iff \cos\theta = 0
\iff$ de richtingen staan loodrecht op elkaar. De nulvector is \omterm{def:g11:scal:orthogonal}{orthogonaal}
met alles.
\end{remark}

\begin{example}
$\vec u\,(3, 2)$ en $\vec v\,(-4, 6)$: $\vec u \cdot \vec v = -12 + 12 = 0$:
\omterm{def:g11:scal:orthogonal}{orthogonaal}. Algemeen is $\vec u\,(a, b)$ altijd \omterm{def:g11:scal:orthogonal}{orthogonaal} met
$\vec v\,(-b, a)$ --- een kwartdraai.
\end{example}

\section{De cosinusregel}

\begin{theorem}[Cosinusregel]\label{thm:g11:scal:cosines}
\index{cosinusregel}
In een driehoek $ABC$ met zijden $a = BC$, $b = CA$, $c = AB$ en hoek
$\widehat A$ in het hoekpunt $A$ geldt
\[
a^2 = b^2 + c^2 - 2bc\,\cos\widehat A .
\]
\end{theorem}

\begin{proof}
Schrijf $\vect{BC} = \vect{BA} + \vect{AC} = \vect{AC} - \vect{AB}$ en werk
uit met \cref{prop:g11:scal:rules}:
\[
a^2 = \norm{\vect{BC}}^2
= \norm{\vect{AC}}^2 - 2\,\vect{AC}\cdot\vect{AB} + \norm{\vect{AB}}^2
= b^2 + c^2 - 2\,\vect{AB}\cdot\vect{AC}.
\]
Volgens de cosinusformule is $\vect{AB}\cdot\vect{AC} = c\,b\,\cos\widehat A$.
\end{proof}

\begin{remark}
Is $\widehat A$ een rechte hoek, dan is $\cos\widehat A = 0$ en herleidt de
cosinusregel zich tot de stelling van Pythagoras $a^2 = b^2 + c^2$: het is
Pythagoras met een correctieterm voor niet-rechte hoeken.
\end{remark}

\begin{omfigure}
\begin{tikzpicture}[scale=1.15]
  \coordinate (A) at (0,0);
  \coordinate (B) at (4,0);
  \coordinate (C) at (1.2,2.2);
  \draw[omDef, thick] (A) -- (B) -- (C) -- cycle;
  \node[below left, font=\small] at (A) {$A$};
  \node[below right, font=\small] at (B) {$B$};
  \node[above, font=\small] at (C) {$C$};
  \node[below, font=\small] at (2,0) {$c$};
  \node[above left, font=\small] at (0.6,1.1) {$b$};
  \node[above right, font=\small] at (2.6,1.1) {$a$};
  \draw[->,omThm,thin] (0.7,0) arc (0:61.4:0.7);
  \node[omThm, font=\small] at (30:1.0) {$\widehat A$};
\end{tikzpicture}

{\small De cosinusregel: twee zijden $b$, $c$ en de hoek ertussen leggen de
derde zijde $a$ vast.}
\end{omfigure}

\begin{example}\label{ex:g11:scal:cosines}
Twee zijden van een driehoek meten $b = 5$ en $c = 8$ en sluiten een hoek
$\widehat A = \frac\pi3$ in. Dan is
\[
a^2 = 25 + 64 - 2 \times 5 \times 8 \times \frac12 = 89 - 40 = 49,
\]
dus $a = 7$.
\end{example}

\section{Toepassingen: cirkels en loodrechte rechten}

\begin{proposition}[Cirkel met gegeven middellijn]\label{prop:g11:scal:diameter}
Een punt $M$ ligt op de cirkel met middellijn $[AB]$ dan en slechts dan als
\[
\vect{MA} \cdot \vect{MB} = 0 .
\]
\end{proposition}

\begin{proof}
Is $M \neq A, B$, dan zegt de voorwaarde dat de hoek in $M$ in de driehoek
$AMB$ recht is, en de punten die een lijnstuk onder een rechte hoek zien zijn
precies die van de cirkel met dat lijnstuk als middellijn. Hier is een
rechtstreeks bewijs: zij $O$ het \omterm{prop:g10:coordgeom:midpoint}{midden} van $[AB]$ en $r = \frac{AB}{2}$.
Dan is $\vect{MA} = \vect{MO} + \vect{OA}$ en
$\vect{MB} = \vect{MO} + \vect{OB} = \vect{MO} - \vect{OA}$, dus
\[
\vect{MA}\cdot\vect{MB}
= \norm{\vect{MO}}^2 - \norm{\vect{OA}}^2 = MO^2 - r^2,
\]
wat precies verdwijnt wanneer $MO = r$, dus wanneer $M$ op de cirkel met
middelpunt $O$ en straal $r$ ligt.
\end{proof}

\begin{proposition}[Normaalvector van een rechte]\label{prop:g11:scal:normal}
De \omterm{def:g11:vect:vector}{vector} $\vec n\,(a, b)$ is \omterm{def:g11:scal:orthogonal}{orthogonaal} met elke \omterm{def:g11:vect:direction}{richtingsvector} van de
rechte $d : ax + by + c = 0$; men noemt $\vec n$ een
\emph{normaalvector}\index{normaalvector} van $d$. Twee rechten staan
loodrecht op elkaar precies wanneer hun normaalvectoren (of hun
\omterm{def:g11:vect:direction}{richtingsvectoren}) \omterm{def:g11:scal:orthogonal}{orthogonaal} zijn.
\end{proposition}

\begin{proof}
Een \omterm{def:g11:vect:direction}{richtingsvector} van $d$ is $\vec u\,(-b, a)$
(\cref{thm:g11:vect:cartesian}), en $\vec n \cdot \vec u = a(-b) + b(a) = 0$.
\end{proof}

\begin{method}[Werken met normaalvectoren]\label{met:g11:scal:normal}
\begin{itemize}
  \item Rechte door $A$ met normaalvector $\vec n\,(a,b)$: schrijf
        $\vect{AM} \cdot \vec n = 0$, wat uitwerkt tot $ax + by + c = 0$,
        waarbij $A$ de waarde van $c$ bepaalt.
  \item Loodrechte stand van $d: ax+by+c = 0$ en $d': a'x+b'y+c' = 0$: ga
        $aa' + bb' = 0$ na.
\end{itemize}
\end{method}

\begin{example}\label{ex:g11:scal:perpline}
De rechte door $P(2, -1)$ loodrecht op $d : 3x - y + 4 = 0$: een
\omterm{def:g11:vect:direction}{richtingsvector} van $d$ is $(1, 3)$, en die dient als \emph{normaalvector}
van de loodlijn. \omterm{def:g10:algebra:equation}{Vergelijking}: $x + 3y + c = 0$ met $2 - 3 + c = 0$, dus
$c = 1$: de rechte is $x + 3y + 1 = 0$.
\end{example}

\begin{example}[Vergelijking van een cirkel]\label{ex:g11:scal:circle}
De cirkel met middelpunt $\Omega(1, 2)$ en straal $3$ is de verzameling van
de punten $M(x,y)$ met $\Omega M^2 = 9$:
\[
(x - 1)^2 + (y - 2)^2 = 9 .
\]
Omgekeerd herschrijft $x^2 + y^2 - 2x - 4y - 4 = 0$ zich, door de kwadraten
aan te vullen (\cref{ch:g11:quad}), tot $(x-1)^2 + (y-2)^2 = 9$: dezelfde
cirkel.
\end{example}

\section{Oefeningen}

\begin{exercise}[$\star$]\label{exo:g11:scal:1}
Bereken $\vec u \cdot \vec v$ voor
\[
\vec u\,(2, 5),\ \vec v\,(3, -1); \qquad
\vec u\,(1, -4),\ \vec v\,(8, 2); \qquad
\norm{\vec u} = 3,\ \norm{\vec v} = 4,\ \theta = \frac{\pi}{3} .
\]
\end{exercise}

\begin{exercise}[$\star$]\label{exo:g11:scal:2}
$ABC$ is een gelijkzijdige driehoek met zijde $6$. Bereken
$\vect{AB} \cdot \vect{AC}$ en $\vect{AB} \cdot \vect{BC}$.
\end{exercise}

\begin{exercise}[$\star$]\label{exo:g11:scal:3}
Voor welke waarde van $t$ zijn $\vec u\,(3, t)$ en $\vec v\,(t - 8, 1)$
\omterm{def:g11:scal:orthogonal}{orthogonaal}? En voor welke waarde of waarden van $t$ is $\vec u\,(t, 2)$
\omterm{def:g11:scal:orthogonal}{orthogonaal} met zichzelf min $\vec v\,(4, 2)$, dus
$\vec u \perp (\vec u - \vec v)$?
\end{exercise}

\begin{exercise}[$\star$]\label{exo:g11:scal:4}
Bereken de hoek $\theta$ tussen $\vec u\,(1, 2)$ en $\vec v\,(3, 1)$ (geef
$\cos\theta$ exact, en daarna $\theta$ op de graad nauwkeurig).
\end{exercise}

\begin{exercise}[$\star\star$]\label{exo:g11:scal:5}
Een driehoek heeft zijden $b = 4$ en $c = 6$ met daartussen de hoek
$\widehat A = \frac{2\pi}{3}$. Bereken de derde zijde $a$. Bereken daarna in
een driehoek met zijden $a = 7$, $b = 5$, $c = 4$ de waarde van
$\cos\widehat A$ en beslis of $\widehat A$ scherp of stomp is.
\end{exercise}

\begin{exercise}[$\star\star$]\label{exo:g11:scal:6}
Zij $A(1, 1)$, $B(5, -1)$ en $C(3, 5)$. Bereken
$\vect{AB} \cdot \vect{AC}$; is de hoek $\widehat A$ recht? Bereken de drie
zijdelengten en toets je antwoord met de cosinusregel.
\end{exercise}

\begin{exercise}[$\star\star$]\label{exo:g11:scal:7}
Geef een \omterm{def:g10:algebra:equation}{vergelijking} van de cirkel met middelpunt $\Omega(-2, 3)$ en straal
$5$. Bepaal daarna het middelpunt en de straal van de cirkel
\[
x^2 + y^2 - 6x + 2y - 6 = 0 .
\]
\end{exercise}

\begin{exercise}[$\star\star$]\label{exo:g11:scal:8}
Zij $A(-3, 0)$ en $B(0, 2)$. Zoek met \cref{prop:g11:scal:diameter} een
\omterm{def:g10:algebra:equation}{vergelijking} van de cirkel met middellijn $[AB]$, en ga na dat de oorsprong
$O(0,0)$ erop ligt.
\end{exercise}

\begin{exercise}[$\star\star$]\label{exo:g11:scal:9}
Zoek een \omterm{def:g10:algebra:equation}{vergelijking} van de rechte door $P(1, 3)$ loodrecht op
$d : 2x + y - 7 = 0$, en de \omterm{def:g10:coordgeom:system}{coördinaten} van de voet $H$ van de loodlijn (het
snijpunt van de twee rechten). Leid de afstand van $P$ tot $d$ af.
\end{exercise}

\begin{exercise}[$\star\star$]\label{exo:g11:scal:10}
Bewijs met $\norm{\vec u + \vec v}^2 = \norm{\vec u}^2 + 2\vec u\cdot\vec v +
\norm{\vec v}^2$ en het analogon voor $\vec u - \vec v$ de
\emph{parallellogramidentiteit}:
\[
\norm{\vec u + \vec v}^2 + \norm{\vec u - \vec v}^2
= 2\norm{\vec u}^2 + 2\norm{\vec v}^2 ,
\]
en duid ze in een parallellogram (diagonalen tegenover zijden).
\end{exercise}

\begin{exercise}[$\star\star\star$]\label{exo:g11:scal:11}
Zij $A$ en $B$ twee punten met $AB = 4$, en $I$ het \omterm{prop:g10:coordgeom:midpoint}{midden} van $[AB]$.
\begin{enumerate}
  \item Toon aan dat voor elk punt $M$ geldt:
        $\vect{MA}\cdot\vect{MB} = MI^2 - 4$. (Schuif $I$ ertussen:
        $\vect{MA} = \vect{MI} + \vect{IA}$.)
  \item Leid de verzameling af van de punten $M$ met
        $\vect{MA}\cdot\vect{MB} = 12$.
\end{enumerate}
\end{exercise}

\section{Opgave: twee geschenken van het scalair product}

\begin{problem}[{Weekendopgave --- de somformules vallen uit één scalair
product, de afstand tot een rechte uit een ander, en de drie hoogtelijnen
ontmoeten elkaar als toegift}]
\label{pb:g11:scal:1}
Het \omterm{def:g11:scal:dot}{scalair product} lijkt boekhouding --- vermenigvuldig, tel op --- en toch
levert het bijna gratis twee schatten op waaraan eeuwen hoekenjagen zich
stukbeten: de \emph{somformules} van de goniometrie (bereken $\cos 15^\circ$
exact!) en de \emph{afstand van een punt tot een rechte} in één nette breuk.
Als afscheidscadeau bewijst het dat de drie hoogtelijnen van elke driehoek
door één punt gaan. Alle vier de gezichten van \cref{thm:g11:scal:faces}
melden zich voor dienst.

\medskip
\noindent\textbf{Deel I --- Vlot rekenen.}
\begin{enumerate}
  \item Bereken voor $\vec u(3, 4)$ en $\vec v(-1, 2)$ het \omterm{def:g11:scal:dot}{scalair product}
        $\vec u \cdot \vec v$, de twee \omterm{def:g11:scal:dot}{normen}, en de hoek tussen de \omterm{def:g11:vect:vector}{vectoren}
        (op een tiende graad).
  \item Zoek $k$ zodat $(2, k)$ en $(3, -4)$ \omterm{def:g11:scal:orthogonal}{orthogonaal} zijn
        (\cref{def:g11:scal:orthogonal}).
  \item Het gezicht van de natuurkunde: een slee wordt $10$ m ver getrokken
        door een kracht van $50$ N onder $60^\circ$ met de grond. Bereken de
        arbeid $W = \vec F \cdot \vec d$.
  \item Cosinusregel (\cref{thm:g11:scal:cosines}): twee zijden van $5$ en
        $7$ sluiten een hoek van $60^\circ$ in; zoek de derde zijde exact.
  \item Omgekeerd: een driehoek heeft zijden $4$, $6$, $8$. Bereken de
        \omterm{def:g11:trigo:cossin}{cosinus} van de hoek tegenover de zijde $8$, bepaal het soort hoek, en
        geef hem op een tiende graad.
\end{enumerate}

\medskip
\noindent\textbf{Deel II --- Eerste geschenk: de somformules.}
Zij $\vec u = (\cos a, \sin a)$ en $\vec v = (\cos b, \sin b)$: twee \omterm{def:g11:vect:vector}{vectoren}
van lengte $1$ onder de hoeken $a$ en $b$.
\begin{enumerate}[resume]
  \item Bereken $\vec u \cdot \vec v$ twee keer --- één keer met \omterm{def:g10:coordgeom:system}{coördinaten},
        één keer met de \omterm{def:g11:scal:dot}{normen} en de hoek tussen de \omterm{def:g11:vect:vector}{vectoren} --- en besluit
        tot de moederformule:
        \[
        \cos(a - b) = \cos a \cos b + \sin a \sin b .
        \]
  \item Vervang $b$ door $-b$ (denk aan de pariteit van \omterm{def:g11:trigo:cossin}{sinus} en \omterm{def:g11:trigo:cossin}{cosinus}!) om
        $\cos(a + b)$ te krijgen.
  \item Verkrijg $\sin(a + b)$ uit de complementaire identiteit
        $\sin x = \cos\left(\frac\pi2 - x\right)$, toegepast op $x = a + b$.
  \item Eindelijk exacte waarden: bereken $\cos 15^\circ$ (als
        $\cos(45^\circ - 30^\circ)$) en $\sin 75^\circ$. Ga beide numeriek
        na.
  \item Stel $b = a$ om de formules voor de dubbele hoek af te leiden:
        $\cos 2a$ (in haar drie vormen) en $\sin 2a$.
  \item Bereken uit $\cos 2a = 1 - 2\sin^2 a$ de waarde van
        $\sin^2 15^\circ$, en daarna $\sin 15^\circ$ als een geneste \omterm{def:g11:quad:discriminant}{wortel}
        --- en ga na dat haar kwadraat overeenstemt met de waarde
        $\frac{\sqrt6 - \sqrt2}{4}$ die de methode van vraag 9 oplevert.
\end{enumerate}

\medskip
\noindent\textbf{Deel III --- Tweede geschenk: de afstand tot een rechte.}
\begin{enumerate}[resume]
  \item Zij $d$ de rechte $ax + by + c = 0$, met normaalvector $\vec n(a, b)$
        (\cref{prop:g11:scal:normal}), en $P(x_0, y_0)$ een punt. Leg uit
        waarom, met een willekeurig punt $A$ op $d$, de afstand van $P$ tot
        $d$ gelijk is aan
        $\dfrac{\abs{\vect{AP} \cdot \vec n}}{\norm{\vec n}}$, en leid de
        formule af
        \[
        \operatorname{dist}(P, d)
        = \frac{\abs{a x_0 + b y_0 + c}}{\sqrt{a^2 + b^2}} .
        \]
        Bereken de afstand van $P(5, 6)$ tot $3x + 4y - 12 = 0$.
  \item Schrijf de \omterm{def:g10:algebra:equation}{vergelijking} op van de cirkel met middelpunt $(5, 6)$ die
        die rechte raakt.
  \item Toets de formule: bereken de afstand van de oorsprong tot
        $3x + 4y - 12 = 0$, en ga ze na door de voet van de loodlijn uit de
        oorsprong expliciet te berekenen.
  \item Driehoek $A(0,0)$, $B(6,0)$, $C(2,5)$: bereken zijn oppervlakte met
        $[AB]$ als basis; bereken daarna de afstand van $B$ tot de rechte
        $(AC)$ en ga na dat ze dezelfde oppervlakte oplevert.
  \item De cirkel met middelpunt $(2, 1)$ en straal $2$, tegenover de rechte
        $x + y - 6 = 0$: bereken de afstand van het middelpunt tot de rechte
        en bepaal hun onderlinge ligging (snijdend, rakend, of zonder
        gemeenschappelijk punt).
\end{enumerate}

\medskip
\noindent\textbf{Deel IV --- De toegift: de hoogtelijnen gaan door één punt.}
\begin{enumerate}[resume]
  \item Veralgemeen \cref{exo:g11:scal:11}: bewijs voor elke twee punten $A$
        en $B$ met \omterm{prop:g10:coordgeom:midpoint}{midden} $I$ dat
        \[
        \vect{MA} \cdot \vect{MB} = MI^2 - \frac{AB^2}{4}
        \quad\text{voor elk punt } M .
        \]
  \item Leid \cref{prop:g11:scal:diameter} in één regel af: de verzameling
        van de punten $M$ met $\vect{MA} \cdot \vect{MB} = 0$ is de cirkel
        met middellijn $[AB]$.
  \item De hoogtelijnen: zij $H$ het snijpunt van de hoogtelijnen uit $A$ en
        uit $B$ in een driehoek $ABC$ (dus $\vect{HA} \cdot \vect{BC} = 0$ en
        $\vect{HB} \cdot \vect{CA} = 0$). Bewijs de identiteit
        \[
        \vect{HA} \cdot \vect{BC} + \vect{HB} \cdot \vect{CA}
        + \vect{HC} \cdot \vect{AB} = 0
        \]
        --- werk alles vanuit $H$ uit met Chasles --- en besluit dat $H$ ook
        op de hoogtelijn uit $C$ ligt: de drie hoogtelijnen gaan door één
        punt (het \emph{hoogtepunt}\index{hoogtepunt}).
  \item Slotstuk: vier gezichten, drie trofeeën. Welk gezicht van
        \cref{thm:g11:scal:faces} droeg de somformules, welk de
        afstandsformule, welk het \omterm{pb:g11:scal:1}{hoogtepunt} --- en waarom is ``bereken
        dezelfde grootheid twee keer'' het motto van het hele hoofdstuk?
\end{enumerate}
\end{problem}
